\documentclass[10pt,a4paper]{amsart}
\usepackage[margin=19mm]{geometry}
\usepackage{amsmath,amssymb,mathtools}
\usepackage[scr=rsfs,cal=euler]{mathalfa}
\usepackage{xfrac}
\usepackage[unicode,hidelinks,bookmarks=false]{hyperref}
\newcommand{\ie}{i.e.\ }
\newcommand{\cf}{cf.\ }
\newcommand{\resp}{respectively\ }
\newcommand{\Subsection}{\subsection}
\providecommand{\nobreakdash}{\nobreak\hbox{-}}
\setlength{\emergencystretch}{2em}
\multlinegap0pt
\usepackage{bm}
\usepackage{bbm}
\usepackage{dsfont}
\usepackage{xspace}
\usepackage{cancel}
\usepackage{mathtools}
\usepackage{amsthm}
\makeatletter
\@ifundefined{theorem}{\newtheorem{theorem}{Theorem}}{}
\@ifundefined{definition}{\newtheorem{definition}[theorem]{Definition}}{}
\@ifundefined{lemma}{\newtheorem{lemma}[theorem]{Lemma}}{}
\makeatother
\newcommand{\N}{\mathbb{N}}
\newcommand{\R}{\mathbb{R}}
\newcommand{\ii}{\mathrm{i}}
\title[Cobham's theorem for the Gaussian integers]{Cobham's theorem for the Gaussian integers}
\author{\'Alvaro Bustos-Gajardo, Robbert Fokkink,, Reem Yassawi}
\date{Source: arXiv:2510.01440v3 (CC BY 4.0)}
\begin{document}
\maketitle

\noindent\emph{Source and licence.} Cobham's theorem for the Gaussian integers. Version arXiv:2510.01440v3 (CC BY 4.0), distributed under the Creative Commons Attribution 4.0 International licence, as recorded in the arXiv licence metadata for this exact version and shown on the abstract page https://arxiv.org/abs/2510.01440v3. Full rights notice: licenses/arxiv-2510.01440-CC-BY-4.0.txt.\par\smallskip

\noindent\textit{Editorial scope.} Counted unit: the Lemma whose source label is \texttt{lem:non-collinearity-1} in the article (source0.tex lines 548--552, source label \texttt{lem:non-collinearity-1}), delivered together with the article's own complete original proof of it (source0.tex lines 555--581, one \texttt{proof} environment of 27 lines). No mathematics is written, repaired or re-derived: statement and proof are the article's own text line for line. Required context retained uncounted: def:return-number (lines 538--543). Every definition, display and label of those lines is retained unchanged. Omitted from this package and not redistributed: 1--537, 544--547, 553--554, 582--1022. Nothing inside a retained excerpt is omitted or altered: every retained body is delivered exactly as the article's lines are written, so no omission receipt of any kind accompanies this package. Citations: the retained excerpts carry no citation command at all, so no bibliography entry of the article is transcribed and no reference is redistributed. Reference adaptation: none was needed and none was made; every reference inside the delivered unit resolves inside it, so no unresolved reference or citation is produced. Printed slips kept literal: no printed word, symbol, hypothesis or number of the counted statement or of its proof was altered. Layout and class: the article's own class is \texttt{amsart} with packages including fontenc, graphicx, amsmath, amsfonts, amssymb, amsthm, tikz, tikz-cd, hyperref, placeins, bm, thmtools, mathtools, mathptmx; the wrapper is the lane's pinned \texttt{amsart} editorial wrapper at 10pt on A4, which restates the theorem-like environments the retained text needs and adds the article's own macro definitions that the retained text needs (\texttt{\textbackslash N}, \texttt{\textbackslash R}, \texttt{\textbackslash ii}), with extra wrapper packages (bm, bbm, dsfont, xspace, cancel, mathtools). The delivered numbering is the unit's own. Assets: no retained span contains a figure environment, a graphics-inclusion command, a PDF-page inclusion, a table or a TikZ picture; no image file is copied into this package, the package declares no asset dependency and no image file is redistributed with it. Licence: version arXiv:2510.01440v3 of this article is distributed under the Creative Commons Attribution 4.0 International licence, as recorded in the arXiv licence metadata for this exact version and shown on the abstract page https://arxiv.org/abs/2510.01440v3. A full rights notice is delivered at licenses/arxiv-2510.01440-CC-BY-4.0.txt. Characters: every retained excerpt is pure ASCII, so the delivered engine, XeLaTeX with its default Latin Modern text font, reports no missing character.
   \begin{definition}\label{def:return-number} Let $\mathcal A_\beta $ be a DFAO, and  let $s\in S_\infty$. Then, by the pigeonhole principle, there exist $s'\in S_\infty$ and three words $w_1,w_2,w_3\in D_\beta^*$ for which
    $\delta_\beta(s_0,w_1)=s'$,  $\delta_\beta(s',w_2)=s'$ and $\delta_\beta(s',w_3)=s$. 
    Furthermore, since $s'\in S_\infty$, there are infinitely many choices for  $[w_1]_\beta$, so we pick one satisfying   $[w_1]_\beta \cdot (1-\beta^{\lvert w_2\rvert}) \ne  [w_2]_\beta$. Note that for each $j$, $u_j=[w_1(w_2)^jw_3]_\beta$ is a number
      whose $(\beta, D_\beta)$-expansion ends at state $s$. 
    We  call such numbers 
      $u_j$ {\em return numbers to $s$}.\end{definition}

\begin{lemma}\label{lem:non-collinearity-1}
    Let $s\in S_\infty$, and  suppose that $(u_j)$  are  return numbers to $s$.
    % \tcr{Define the Gaussian integers $u_j=[w_1w_2^j]_\beta\in\Z[\ii]$.} \tcr{If for some $k\ge 0,\ell>0$ the Gaussian integers $u_k,u_{k+\ell},u_{k+2\ell}$ are collinear} 
      If there exists an  arithmetic progression of indices such that  $u_k,u_{k+\ell},u_{k+2\ell}$ are collinear return numbers, then $\beta$ is the root of an integer.
\end{lemma}

\begin{proof}

We have $u_j=[w_1w_2^j w_3]_\beta$ where the words $w_i$ are as in Definition~\ref{def:return-number}. First suppose that $s=s'$ and $w_3$ is the empty word, i.e., that there is a cycle rooted at $s$.

    Let $N=\lvert w_2\rvert$ and $h=[w_2]_\beta$. The sequence of return numbers $(u_j)_{j\in\N}$  satisfies the recurrence relation $u_{j+1} = \beta^N u_j + h$, from which the following equalities immediately follow:
        \begin{align*}
            u_{k+\ell} &= \beta^{N\ell}u_k+h(\beta^{N(\ell - 1)} + \beta^{N(\ell-2)} + \cdots + \beta ^N + 1)  \\
            u_{k+2\ell} &= \beta^{N\ell}u_{k+\ell}+h(\beta^{N(\ell - 1)} + \beta^{N(\ell-2)} + \cdots + \beta ^N + 1)  \\
            {}\implies u_{k+2\ell}-u_{k+\ell} &= \beta^{N\ell}(u_{k+\ell}-u_k).
        \end{align*}
        
    We now face two scenarios. If $u_{k+\ell}-u_k\ne 0$, we may divide by this term and obtain the equality $\beta^{N\ell}=\frac{u_{k+2\ell}-u_{k+\ell}}{u_{k+\ell}-u_k}$. If $u_k,u_{k+\ell},u_{k+2\ell}$ lie on the same line in the plane, then $u_{k+2\ell}-u_{k+\ell}$ and $u_{k+\ell}-u_{k}$ must lie on a parallel line which passes through the origin, i.e. there exists a real constant $\lambda$ for which $u_{k+2\ell}-u_{k+\ell} = \lambda(u_{k+\ell}-u_{k})$. But then $\beta^{N\ell}=\lambda\in\R$, which is only possible if $\arg(\beta)$ is a rational multiple of $2\pi\ii$. 
    

    Otherwise, if $u_{k+\ell}-u_k=0$, we obtain that $u_{j+\ell}=u_j$ for all $j\ge k$ from the aforementioned recurrence relation, so the sequence $(u_j)_{j\in\N}$ remains bounded. Write $\hat{u}_j=u_j-\frac{h}{1-\beta^N}$; from this definition, it holds that:
        \begin{align*}
            \hat{u}_{j+1}&=u_{j+1}-\frac{h}{1-\beta^N}\\
            &=\beta^N u_j+h-\frac{h}{1-\beta^N} \\
            &=\beta^N\left(\hat{u}_j+\frac{h}{1-\beta^N}\right)+h\left( 1-\frac{1}{1-\beta^N}\right)\\
            &=\beta^N\hat{u}_j+\frac{\beta^Nh}{1-\beta^N} + \frac{h((1-\beta^N)-1)}{1-\beta^N}\\
            &=\beta^N\hat{u}_j,
        \end{align*}
    and since $\lvert\beta\rvert>1$, we have that $\lvert \hat{u}_j\rvert\to\infty$ unless $\hat{u}_0=u_0-\frac{h}{1-\beta^N}=0$. Since $u_0=[w_1]_\beta$ and $h=[w_2]_\beta$, this is impossible as we chose $w_1$ and $w_2$ so that $[w_1]_\beta \cdot (1-\beta^{\lvert w_2\rvert}) \ne  [w_2]_\beta$. As the sequence $(u_j)_{j\in\N}$ is bounded if and only if $(\hat{u}_j)_{j\in\N}$ remains bounded as well, we get a contradiction, so this scenario never happens.
    
    Consider now the general case when $s\neq s'$.   We define $M=\lvert w_3\rvert$, $h' = [w_3]_\beta$, and for each $j$, $u_j=[w_1(w_2)^j]_\beta$, which are return numbers to $s'$.
    Then we set    $v_k=\beta^M u_k+h'$, which are return numbers to $s$, and  each term $v_k$ is obtained from the term $u_k$ via a dilation followed by a translation. As these transformations preserve collinearity, the  return numbers $v_k, v_{k+\ell}, v_{k+2\ell}$ lie on the same line if and only if the respective  $u_k, u_{k+\ell}, u_{k+2\ell}$ are collinear as well, and thus an application of the first case above to the state $s'$ shows us that $\arg(\beta)$ is a rational multiple of $2\pi$. 
\end{proof}

\end{document}
